\subsection{nml\_output\_variables}
\label{SEC:nml_output_variables}
\inp*[nmlOutputVariables]{}[nml\_output\_variables][]

This \hyperref[SEC:User_Input]{namelist block}
is used to specify which variables to output
to the CGNS files.

\begin{nmlblocklst}[
,label={LST:nml_output_variables}
,caption={Namelist nml\_output\_variables.}
]
&nml_output_variables
  Conserved_Variables            = F
  Momentum                       = F
  Vorticity                      = F
  Cylindrical_Velocity           = F
  Cylindrical_Coordinates        = F
  Relative_Velocity              = F
  Relative_Conditions            = F
  Thermo_Variables               = F
  Temperature                    = F
  Gradients                      = F
  Density_Gradient               = F
  Density_Gradient_Magnitude     = F
  Pressure_Gradient              = F
  Pressure_Gradient_Magnitude    = F
  Pressure_Dmat_Gradient         = F
  Temperature_Gradient           = F
  Temperature_Gradient_Magnitude = F
  Temperature_Dmat_Gradient      = F
  Schlieren                      = F
  Modified_Schlieren             = F
  Transport_Variables            = F
  Stagnation_Variables           = F
  Stagnation_Pressure            = F
  Stagnation_Temperature         = F
  Marker_Variables               = F
  Qcriterion                     = F
  Grid_Details                   = F
  Local_Timestep                 = F
  Species_Mass_Fraction          = T
  Species_Density                = F
  Species_Mole_Fraction          = F
  Extra_Species_Variables        = F
  Debug_Species_Variables        = F
  Grdnt_Species_Variables        = F
  Local_Filter_Dissipation       = F
  CellAve_Density                = F
  Exact_Density                  = F
  Exact_Pressure                 = F
/
\end{nmlblocklst}

%% For future reference:
%%   The default output variable list is in
%%   subroutine get_cgns_sol_variables of write_cgns_mod.f90.

\FINALFIX{\newpage}
The following variables are automatically output to the CGNS files:
\\
  \ov{Density},
  \ov{VelocityX},
  \ov{VelocityY},
  \ov{VelocityZ},
  \ov{Pressure}, and
  \ov{Mach}.

\vspace{1.0ex}
\noindent\textbf{NOTES:}\\
\begin{itemize}
  \item
    The grid coordinates are output in the \key[key:GridUnits]{grid units}
    of the input grid file.  For example, say your grid file has dimensions of
    inches. You will need to include the input
    $\inp[nmlGrid::GridScalingFactor]{Grid\_Scaling\_Factor}=0.0254$ to convert
    the grid file to meters for it to be used correctly in the solver. However,
    this conversion will be undone when creating the CGNS solution files, which
    will contain coordinate dimensions of inches just like the grid file.
  \item
    All variables are output in SI units (except for the coordinate variables as
    discussed in the previous bullet).
  \item
    Due to the first two bullets, if using post-processing software to compute
    integrated quantities, care needs to be taken to ensure the coordinate
    variables and variable(s) being integrated are in compatible units.
  \item
    In the following descriptions, any vector quantities that are output for 1D
    and 2D cases will exclude vector components from higher spatial dimensions.
    For example, \ov{VelocityY} and \ov{VelocityZ} will not be included for a 1D
    flow, and for a gradient variable from a 2D simulation, the
    \ov{[VariableName]GradientZ} variables would not be included.
  \item
    If multiple inputs add the same variable to the list of output variables,
    only one instance of the variable will be written, \thatis any duplicates in
    the list of output variables will be removed.
  \item
    The names of the variables written to the CGNS solution files attempt to
    follow the variable naming conventions of the CGNS standard, \forexample the
    $x$-coordinate variable will have the name \ov{CoordinateX} instead of just
    \quoterm{\textsf{X}}.
\end{itemize}

\devNotes{
CHANGECODE:
Default variables can not be disabled.
Instead of a hard-coded default list, consider adding individual options
to the namelist above with defaults:
\\\indent\raggedright
  \inptext{Density = T},
  \inptext{Velocity = T},
  \inptext{Pressure = T}, and
  \inptext{Mach = T}
}
%==========================================================================

\begin{nmldescitem}
Conserved_Variables = F
\end{nmldescitem}
\inp*[nmlOutputVariables::ConservedVariables]{}[Conserved\_Variables][]

\begin{nmldescbody}
Setting this to true adds the following variables to the CGNS solution files:
\\
  \ov{Density},
  \ov{MomentumX},
  \ov{MomentumY},
  \ov{MomentumZ}, and
  \ov{EnergyStagnationDensity} (\thatis $\rho e_t$).
\end{nmldescbody}

%==========================================================================

\begin{nmldescitem}
Momentum = F
\end{nmldescitem}
\inp*[nmlOutputVariables::Momentum]{}[Momentum][]

\begin{nmldescbody}
Setting this to true adds the following variables to the CGNS solution files:
\\
  \ov{MomentumX},
  \ov{MomentumY}, and
  \ov{MomentumZ}.
\end{nmldescbody}

%==========================================================================

\begin{nmldescitem}
Vorticity = F
\end{nmldescitem}
\inp*[nmlOutputVariables::Vorticity]{}[Vorticity][]

\begin{nmldescbody}
Setting this to true adds the following variables to the CGNS solution files:
\\\raggedright
  \ov{VorticityX},
  \ov{VorticityY},
  \ov{VorticityZ},
  \ov{VorticityMagnitude},
  \ov{DivergenceOfVelocity},
  \ov{StrainRateMagnitude},
  \ov{Qcriterion}, and
  \ov{Qcriterion\_Positive}.
\end{nmldescbody}

%==========================================================================

\begin{nmldescitem}
Cylindrical_Velocity = F
\end{nmldescitem}
\inp*[nmlOutputVariables::CylindricalVelocity]{}[Cylindrical\_Velocity][]

\begin{nmldescbody}
Setting this to true adds the following variables to the CGNS solution files:
\\
  \ov{VelocityR} and
  \ov{VelocityTheta}.
\\
If either
\inp[nmlOutputVariables::RelativeVelocity]{Relative\_Velocity} or
\inp[nmlOutputVariables::RelativeConditions]{Relative\_Conditions}
is also true, then
the following variables will also be added to the CGNS solution files:
\\
  \ov{RelativeVelocityR} and
  \ov{RelativeVelocityTheta}.
\end{nmldescbody}

%==========================================================================

\begin{nmldescitem}
Cylindrical_Coordinates = F
\end{nmldescitem}
\inp*[nmlOutputVariables::CylindricalCoordinates]{}[Cylindrical\_Coordinates][]

\begin{nmldescbody}
Setting this to true adds the following variables to the CGNS \emph{grid} file:
\\
  \ov{CoordinateAxial},
  \ov{CoordinateR}, and
  \ov{CoordinateTheta}.
\\
This will also automatically set 
\inp[nmlOutputVariables::CylindricalVelocity]{Cylindrical\_Velocity} to true.

\noindent\NOTE
The axial coordinate direction will be set to the $x$-axis, \thatis
$\ov{CoordinateAxial}=\ov{CoordinateX}$, unless:
  \begin{itemize}
    \item 
      $\inp[nmlReference::AoAAlpha]{AoA\_Alpha}=\pm90$\textdegree{} and
      $\inp[nmlReference::AoABeta]{AoA\_Beta}=0$\textdegree{}:
      $\quad\Rightarrow\quad \ov{CoordinateAxial}=\ov{CoordinateY}$
    \item 
      $\inptext{AoA\_Alpha}=0$\textdegree{} and
      $\inptext{AoA\_Beta}=\pm90$\textdegree{}:
      $\quad\Rightarrow\quad \ov{CoordinateAxial}=\ov{CoordinateZ}$
  \end{itemize}
The other two cylindrical coordinate variables are determined by the axial
direction and the right-hand rule.

\noindent\NOTE
These variables will only be written to the CGNS grid file. These variables will
be added to each of the solution files by using CGNS links.
\end{nmldescbody}

%==========================================================================

%\FINALFIX{\newpage}
\begin{nmldescitem}
Relative_Velocity = F
\end{nmldescitem}
\inp*[nmlOutputVariables::RelativeVelocity]{}[Relative\_Velocity][]

\begin{nmldescbody}
Setting this to true adds the following variables to the CGNS solution files:
\\
  \ov{RelativeVelocityX},
  \ov{RelativeVelocityY}, and
  \ov{RelativeVelocityZ}.
\end{nmldescbody}

%==========================================================================

\begin{nmldescitem}
Relative_Conditions = F
\end{nmldescitem}
\inp*[nmlOutputVariables::RelativeConditions]{}[Relative\_Conditions][]

\begin{nmldescbody}
Setting this to true adds the following variables to the CGNS solution files:
\\
  \ov{RelativeVelocityX},
  \ov{RelativeVelocityY},
  \ov{RelativeVelocityZ},
\\
  \ov{RelativeMomentumX},
  \ov{RelativeMomentumY},
  \ov{RelativeMomentumZ}, and
  \ov{RelativeMach}.
\end{nmldescbody}

%==========================================================================

\begin{nmldescitem}
Thermo_Variables = F
\end{nmldescitem}
\inp*[nmlOutputVariables::ThermoVariables]{}[Thermo\_Variables][]

\begin{nmldescbody}
Setting this to true adds the following variables to the CGNS solution files:
\\
  \ov{Pressure},
  \ov{Temperature},
  \ov{Enthalpy},
  \ov{SpecificHeatPressure} (\thatis $c_p$), and
  \ov{Entropy}.
\end{nmldescbody}

%==========================================================================

\begin{nmldescitem}
Temperature = F
\end{nmldescitem}
\inp*[nmlOutputVariables::Temperature]{}[Temperature][]

\begin{nmldescbody}
Setting this to true adds the following variable to the CGNS solution files:
  \ov{Temperature}.
\end{nmldescbody}

%==========================================================================

\begin{nmldescitem}
Gradients = F
\end{nmldescitem}
\inp*[nmlOutputVariables::Gradients]{}[Gradients][]

\begin{nmldescbody}
Setting this to true will also include in the CGNS solution file the
directional gradient components for any of the following variables that are
also set to be included in the CGNS solution file:
  \ov{Density},
  \ov{Velocity},
  \ov{EnergyStagnationDensity},
  \ov{Temperature},
  \ov{Momentum},
  \ov{MassFraction}, and
  \ov{SpeciesDensity}.

The name of the gradient variables will be the name of the
non-gradient variable with
  \ov{GradientX}, \ov{GradientY}, and \ov{GradientZ}
appended on to the variable name,
\forexample
\\
  \ov{MomentumYGradientZ}.
\end{nmldescbody}

%==========================================================================

\begin{nmldescitem}
Density_Gradient = F
\end{nmldescitem}
\inp*[nmlOutputVariables::DensityGradient]{}[Density\_Gradient][]

\begin{nmldescbody}
Setting this to true will add the directional gradient components of density
to the CGNS solution files, with the names of the variables being
\ov{DensityGradientX}, \ov{DensityGradientY}, and \ov{DensityGradientZ}.
\end{nmldescbody}

%==========================================================================

\FINALFIX{\newpage}
\begin{nmldescitem}
Density_Gradient_Magnitude = F
\end{nmldescitem}
\inp*[nmlOutputVariables::DensityGradientMagnitude]{}[Density\_Gradient\_Magnitude][]

\begin{nmldescbody}
Setting this to true will add the magnitude of the density gradient
to the CGNS solution files under the variable name \ov{GradDensityMagnitude}.
\end{nmldescbody}

%==========================================================================

\begin{nmldescitem}
Pressure_Gradient = F
\end{nmldescitem}
\inp*[nmlOutputVariables::PressureGradient]{}[Pressure\_Gradient][]

\begin{nmldescbody}
Setting this to true will add the directional gradient components of pressure
to the CGNS solution files, with the names of the variables being
\ov{PressureGradientX}, \ov{PressureGradientY}, and \ov{PressureGradientZ}.

\noindent\NOTE
This will use the chain rule to compute a pressure gradient that is consistent
with the gradients of the conserved variables.
\end{nmldescbody}

%==========================================================================

\begin{nmldescitem}
Pressure_Gradient_Magnitude = F
\end{nmldescitem}
\inp*[nmlOutputVariables::PressureGradientMagnitude]{}[Pressure\_Gradient\_Magnitude][]

\begin{nmldescbody}
Setting this to true will add the magnitude of the pressure gradient
to the CGNS solution files under the variable name \ov{GradPresMagnitude}.

\noindent\NOTE
This will use the chain rule to compute a pressure gradient that is consistent
with the gradients of the conserved variables.
\end{nmldescbody}

%==========================================================================

\begin{nmldescitem}
Pressure_Dmat_Gradient = F
\end{nmldescitem}
\inp*[nmlOutputVariables::PressureDmatGradient]{}[Pressure\_Dmat\_Gradient][]

\begin{nmldescbody}
Setting this to true will use the derivative matrices to compute the gradient
of the pressure and add the resulting directional gradient components and the
gradient magnitude to the CGNS solution files, with the names of the variables
being \ov{PressureDmatrixX}, \ov{PressureDmatrixY}, \ov{PressureDmatrixZ}, and
\ov{PressureDmatrixMagnitude}.

\noindent\NOTE
The pressure field in a cell is not in the solution polynomial space, so the
resulting pressure gradients will most likely contain some amount of aliasing
errors since derivative matrices are only exact for a function within the
solution polynomial space.
\end{nmldescbody}

%==========================================================================

\begin{nmldescitem}
Temperature_Gradient = F
\end{nmldescitem}
\inp*[nmlOutputVariables::TemperatureGradient]{}[Temperature\_Gradient][]

\begin{nmldescbody}
Setting this to true will add the directional gradient components of temperature
to the CGNS solution files, with the names of the variables being
\ov{TemperatureGradientX}, \ov{TemperatureGradientY}, and \ov{TemperatureGradientZ}.

\noindent\NOTE
This will use the chain rule to compute a temperature gradient that is consistent
with the gradients of the conserved variables.
\end{nmldescbody}

%==========================================================================

\begin{nmldescitem}
Temperature_Gradient_Magnitude = F
\end{nmldescitem}
\inp*[nmlOutputVariables::TemperatureGradientMagnitude]{}[Temperature\_Gradient\_Magnitude][]

\begin{nmldescbody}
Setting this to true will add the magnitude of the temperature gradient
to the CGNS solution files under the variable name \ov{GradTempMagnitude}.

\noindent\NOTE
This will use the chain rule to compute a temperature gradient that is consistent
with the gradients of the conserved variables.
\end{nmldescbody}

%==========================================================================

\FINALFIX{\newpage}
\begin{nmldescitem}
Temperature_Dmat_Gradient = F
\end{nmldescitem}
\inp*[nmlOutputVariables::TemperatureDmatGradient]{}[Temperature\_Dmat\_Gradient][]

\begin{nmldescbody}
Setting this to true will use the derivative matrices to compute the gradient of
the temperature and add the resulting directional gradient components and the
gradient magnitude to the CGNS solution files, with the names of the variables
being \ov{TemperatureDmatrixX}, \ov{TemperatureDmatrixY},
\ov{TemperatureDmatrixZ}, and \ov{TemperatureDmatrixMagnitude}.

\noindent\NOTE
The temperature field in a cell is not in the solution polynomial space, so the
resulting temperature gradients will most likely contain some amount of aliasing
errors since derivative matrices are only exact for a function within the
solution polynomial space.
\end{nmldescbody}

%==========================================================================

\begin{nmldescitem}
Schlieren = F
\end{nmldescitem}
\inp*[nmlOutputVariables::Schlieren]{}[Schlieren][]

\begin{nmldescbody}
Setting this to true will compute a numerical Schlieren variable to include
in the CGNS solution files under the variable name \ov{Schlieren}.

\noindent\NOTE
The numerical Schlieren is simply the magnitude of the density gradient, \thatis
$\nabla\rho$.
%
This is not to be confused with what is sometimes referred to as
``shadowgraph'' which is the Laplacian of density, \thatis $\nabla^2\rho$.
\end{nmldescbody}

%==========================================================================

\begin{nmldescitem}
Modified_Schlieren = F
\end{nmldescitem}
\inp*[nmlOutputVariables::ModifiedSchlieren]{}[Modified\_Schlieren][]

\begin{nmldescbody}
Setting this to true will add the variable \ov{ModifiedSchlieren} to the
CGNS solution files, which is a $\log_{10}$ scaling of the numerical Schlieren
variable, \thatis $\log_{10}\left(1+\nabla\rho\right)$ where the offset by $1$
is simply to avoid $\log_{10}(0)$ when the density gradient is $0$.
\end{nmldescbody}

%==========================================================================

\begin{nmldescitem}
Transport_Variables = F
\end{nmldescitem}
\inp*[nmlOutputVariables::TransportVariables]{}[Transport\_Variables][]

\begin{nmldescbody}
Setting this to true adds the following variables to the CGNS solution files:
  \ov{ViscosityMolecular} and
  \ov{ThermalConductivity}.
\\
If using a multi-component gas mixture and
\inp[nmlOutputVariables::ExtraSpeciesVariables]{Extra\_Species\_Variables}
is also true, then the variable
  \ov{Prandtl}
giving the mixture Prandtl number will also be included in the
CGNS solution file.
\end{nmldescbody}

%==========================================================================

\begin{nmldescitem}
Stagnation_Variables = F
\end{nmldescitem}
\inp*[nmlOutputVariables::StagnationVariables]{}[Stagnation\_Variables][]

\begin{nmldescbody}
Setting this to true adds the following variables to the CGNS solution files:
  \ov{PressureStagnation},
  \ov{TemperatureStagnation}, and
  \ov{EnthalpyStagnation}.
\end{nmldescbody}

%==========================================================================

\begin{nmldescitem}
Stagnation_Pressure = F
\end{nmldescitem}
\inp*[nmlOutputVariables::StagnationPressure]{}[Stagnation\_Pressure][]

\begin{nmldescbody}
Setting this to true adds the following variable to the CGNS solution files:
  \ov{PressureStagnation}.
\end{nmldescbody}

%==========================================================================

\begin{nmldescitem}
Stagnation_Temperature = F
\end{nmldescitem}
\inp*[nmlOutputVariables::StagnationTemperature]{}[Stagnation\_Temperature][]

\begin{nmldescbody}
Setting this to true adds the following variable to the CGNS solution files:
  \ov{TemperatureStagnation}.
\end{nmldescbody}

%==========================================================================

\FINALFIX{\newpage}
\begin{nmldescitem}
Marker_Variables = F
\end{nmldescitem}
\inp*[nmlOutputVariables::MarkerVariables]{}[Marker\_Variables][]

\begin{nmldescbody}
If using \inp[nmlLimiter::MarkerOption]{Marker\_Option},
the following variables will be added to the CGNS solution files:
\\
%\begin{tabular}{ @{\quad} c @{ \quad } p{0.85\linewidth} @{} }
\begin{tabular}{ @{\quad} c l l l @{} }
  \inptext{Marker\_Option} & \quad & \multicolumn{2}{c}{Output Variables} \\
  \cline{1-1}                         \cline{3-4}
  1 && \texttt{Resolution\_Marker}& \texttt{Scaled\_Resolution\_Marker} \\
  2 && \texttt{Jump\_Marker}      & \texttt{Scaled\_Jump\_Marker} \\
  3 && \texttt{MLP\_Marker}       & \texttt{Scaled\_MLP\_Marker}
\end{tabular}
\\
The variable specified by
  \inp[nmlLimiter::MarkerVariable]{Marker\_Variable}
will also be output.
\end{nmldescbody}

%==========================================================================

\begin{nmldescitem}
Qcriterion = F
\end{nmldescitem}
\inp*[nmlOutputVariables::Qcriterion]{}[Qcriterion][]

\begin{nmldescbody}
Setting this to true adds the following variables to the CGNS solution files:
\\
  \ov{Qcriterion} and
  \ov{Qcriterion\_Positive}.

For the variable \ov{Qcriterion}, the computed Q-criterion is
normalized then rescaled such that the average Q-criterion is 0.0,
and the values $\pm 1$ respectively represent $\pm 1/100$ of a
standard deviation from the average.
%
The variable \ov{Qcriterion\_Positive} contains the true values
of the Q-criterion, but only the positive parts, with the
negative parts set to zero.
\end{nmldescbody}

%==========================================================================

\begin{nmldescitem}
Grid_Details = F
\end{nmldescitem}
\inp*[nmlOutputVariables::GridDetails]{}[Grid\_Details][]

\begin{nmldescbody}
Setting this to true adds the following variables to the CGNS solution files:
\\
  \ov{Partition},
  \ov{CellVolume}, and
  \ov{CellLengthScale}.

The variable \ov{Partition} will contain the partition ID for each grid cell.
%
The variable \ov{CellVolume} will contain the area or volume of each grid
cell depending on whether the grid is 2D or 3D, respectively.
%
The variable \ov{CellLengthScale} will contain an approximate length scale for
each grid cell which is computed as either the square-root of the cell area
for a 2D grid or the cube-root of the cell volume for a 3D grid.
\end{nmldescbody}

%==========================================================================

\begin{nmldescitem}
Local_Timestep = F
\end{nmldescitem}
\inp*[nmlOutputVariables::LocalTimestep]{}[Local\_Timestep][]

\begin{nmldescbody}
Setting this to true adds the following variables to the CGNS solution files:
\\
  \ov{LocalTimestep} and
  \ov{LocalTimestep\_Scaled}.

The variable \ov{LocalTimestep} will contain the physical time step size
computed for each grid cell, whereas the variable \ov{LocalTimestep\_Scaled}
will contain that same value normalized by the minimum time step size across
all grid cells.
\end{nmldescbody}

%==========================================================================

\begin{nmldescitem}
Species_Mass_Fraction = T
\end{nmldescitem}
\inp*[nmlOutputVariables::SpeciesMassFraction]{}[Species\_Mass\_Fraction][]

\begin{nmldescbody}
If using a multi-component gas mixture, setting this to true adds the
mass fractions for each of the gas species to the CGNS solution files.
%
The variable name will be \ov{MassFraction} with the species' chemical
formula appended to the end, \forexample \ov{MassFractionH2O}.
\end{nmldescbody}

%==========================================================================

\begin{nmldescitem}
Species_Density = F
\end{nmldescitem}
\inp*[nmlOutputVariables::SpeciesDensity]{}[Species\_Density][]

\begin{nmldescbody}
If using a multi-component gas mixture, setting this to true adds the
individual density for each of the gas species to the CGNS solution files.
%
The variable name will be \ov{SpeciesDensity} with the species' chemical
formula appended to the end, \forexample \ov{SpeciesDensityH2O}.
\end{nmldescbody}

%==========================================================================

\begin{nmldescitem}
Species_Mole_Fraction = F
\end{nmldescitem}
\inp*[nmlOutputVariables::SpeciesMoleFraction]{}[Species\_Mole\_Fraction][]

\begin{nmldescbody}
If using a multi-component gas mixture, setting this to true adds the
mole fractions for each of the gas species to the CGNS solution files.
%
The variable name will be \ov{MoleFraction} with the species' chemical
formula appended to the end, \forexample \ov{MoleFractionH2O}.
\end{nmldescbody}

%==========================================================================

\begin{nmldescitem}
Extra_Species_Variables = F
\end{nmldescitem}
\inp*[nmlOutputVariables::ExtraSpeciesVariables]{}[Extra\_Species\_Variables][]

\begin{nmldescbody}
If using a multi-component gas mixture, setting this to true adds the
following variables to the CGNS solution files:
\\\raggedright
  \ov{SpeciesDensity},
  \ov{Enthalpy},
  \ov{StagnationEnthalpy},
  \ov{SpecificHeatRatio} (\thatis mixture $\gamma$),
  \ov{SpecificHeatPressure} (\thatis mixture $c_p$),
  \ov{IdealGasConstant} (\thatis mixture $R_{\mathrm{gas}}$),
  \ov{MomentumX},
  \ov{MomentumY},
  \ov{MomentumZ},
  \ov{EnergyStagnationDensity},
  \ov{PressureStagnation}, and
  \ov{TemperatureStagnation}.

Each gas component will have its own \ov{SpeciesDensity} variable
with its respective chemical formula appended to the end,
\forexample \ov{SpeciesDensityCO2}.
\end{nmldescbody}

%==========================================================================

\begin{nmldescitem}
Debug_Species_Variables = F
\end{nmldescitem}
\inp*[nmlOutputVariables::DebugSpeciesVariables]{}[Debug\_Species\_Variables][]

\begin{nmldescbody}
If using a multi-component gas mixture, setting this to true will output the
individual contribution of each species to various mixture thermodynamic
properties, \forexample enthalpy, $c_p$, etc.

\NOTE
These variables are intended to aid developers in debugging the
thermodynamic properties of gas mixtures.
\end{nmldescbody}

%==========================================================================

\begin{nmldescitem}
Grdnt_Species_Variables = F
\end{nmldescitem}
\inp*[nmlOutputVariables::GrdntSpeciesVariables]{}[Grdnt\_Species\_Variables][]

\begin{nmldescbody}
If using a multi-component gas mixture, this is essentially the same as
\inp[nmlOutputVariables::Gradients]{Gradients}, but it also outputs
the gradients of the species mass fractions and densities.
\end{nmldescbody}

%==========================================================================

%\AddEntry[nmlOutputVariables::LocalFilterDissipation]{nml\_output\_Variables}{Local\_Filter\_Dissipation}

\begin{nmldescitem}
Local_Filter_Dissipation = F
\end{nmldescitem}
\inp*[nmlOutputVariables::LocalFilterDissipation]{}[Local\_Filter\_Dissipation][]

\begin{nmldescbody}
This output can be used to assess filter dissipation.
%\vspace*{0.5em}%
%\FIXME{
%\begin{enumerate}
%\item
%What do these variables represent? Units?
%\item
%Proof this description and verify the option is coded.
%\end{enumerate}
%}
\end{nmldescbody}

%==========================================================================

\begin{nmldescitem}
CellAve_Density = F
\end{nmldescitem}
\inp*[nmlOutputVariables::CellAveDensity]{}[CellAve\_Density][]

\begin{nmldescbody}
Setting this to true computes the cell average of the instantaneous density
polynomial to include in the CGNS solution files under the variable name
\ov{CellAveDensity}.
\end{nmldescbody}

%==========================================================================

\begin{nmldescitem}
Exact_Density = F
Exact_Pressure = F
\end{nmldescitem}
\inp*[nmlOutputVariables::ExactDensity]{}[Exact\_Density][]
\inp*[nmlOutputVariables::ExactPressure]{}[Exact\_Pressure][]

\begin{nmldescbody}
For cases that have an analytical solution, the exact density and pressure
can be computed at the output points and written to the solution file
for direct comparison with the numerical solution.
%
Cases that support this output include the following:
\\
\begin{tabular}{cl}
  \inp[nmlInitialization::Initialization]{Initialization} & Description \\
  23 & 1D Sod shock tube \\
  26 & 1D Jiang-Shu density pulse test
\end{tabular}
\end{nmldescbody}



